\chapter{Execution Modes and Artifact Contracts}

The stage entry policy is encoded in \codepath{src/card_framework/runtime/pipeline_plan.py}. The runtime accepts four start stages: \texttt{stage-1}, \texttt{stage-2}, \texttt{stage-3}, and \texttt{stage-4}. Table~\ref{tab:start-stages} summarizes the required inputs and outputs.

\begin{table}[H]
\centering
\small
\begin{tabularx}{\textwidth}{L{1.5cm}L{3.2cm}L{4.4cm}Y}
\toprule
Start stage & Required existing artifact(s) & Executed components & Principal output(s) \\
\midrule
\texttt{stage-1} & input audio & stage-1, stage-2, optional stage-3, optional stage-4 & transcript, summary, optional cloned audio, optional interjected audio \\
\texttt{stage-2} & transcript JSON & stage-2, optional stage-3, optional stage-4 & summary, optional cloned audio, optional interjected audio \\
\texttt{stage-3} & summary XML and transcript metadata that resolves speaker samples & stage-3, optional stage-4 & stage-3 manifest and merged cloned audio, optional stage-4 output \\
\texttt{stage-4} & summary XML plus stage-3 manifest & stage-4 only & interjector manifest and final interjected merged audio \\
\bottomrule
\end{tabularx}
\caption{Execution modes accepted by the runtime.}
\label{tab:start-stages}
\end{table}

The principal serialized contracts are:

\begin{itemize}
  \item \textbf{Transcript JSON:} top-level \codepath{segments} plus \codepath{metadata}, represented at runtime by \codepath{src/card_framework/orchestration/transcript.py}. Extra segment keys are preserved, which is essential for affect metadata and overlap metadata.
  \item \textbf{Summary XML:} speaker-tagged XML turns represented by \codepath{src/card_framework/shared/summary_xml.py}. New outputs are vector-first and can carry \codepath{emotion_vector}, \codepath{emotion_vector_source}, \codepath{source_segment_indices}, and \codepath{source_segment_weights}.
  \item \textbf{Speaker-sample manifest:} generated by \codepath{src/card_framework/audio_pipeline/speaker_samples.py}; it records selected slices, target duration, audio format, and output paths.
  \item \textbf{Voice-clone manifest:} generated by stage-3; it records one synthesized artifact per summary turn and the optional merged WAV.
  \item \textbf{Interjector manifest:} generated by stage-4; it records accepted and rejected overlays, timing source, timing variant, repair history, and overlay-fit telemetry.
\end{itemize}

Listing~\ref{lst:summaryxml} shows the downstream summary shape used by stage-3 and stage-4.

\begin{lstlisting}[style=cfxml,caption={Representative summary XML shape used downstream by stage-3 and stage-4. Adapted from \codepath{src/card_framework/shared/summary_xml.py}.},label={lst:summaryxml}]
<SPEAKER_00 emotion_categories="calm,warm"
            emotion_vector="0.060,0.000,0.020,0.000,0.000,0.090,0.010,0.310"
            emotion_vector_source="source_linked"
            source_segment_indices="17,18,19"
            source_segment_weights="0.410000,0.360000,0.230000">
  The host reframes the episode around recovery, isolation, and renewed purpose.
</SPEAKER_00>
<SPEAKER_01 emotion_categories="surprised,upbeat,calm"
            emotion_vector="0.140,0.000,0.000,0.000,0.000,0.030,0.220,0.180"
            emotion_vector_source="explicit">
  The guest answers with a more vivid, more personal version of the same arc.
</SPEAKER_01>
\end{lstlisting}

\implnote{Transcript DTOs are defined in \codepath{src/card_framework/orchestration/transcript.py}; A2A request/response DTOs are defined in \codepath{src/card_framework/agents/dtos.py}; summary XML parsing and serialization are implemented in \codepath{src/card_framework/shared/summary_xml.py}.}
